\section{Threats to Validity}
The constructs rely on complementary but imperfect proxies. Self-Reported AI
Usage and Perceptions (M1) and Perceived Role Disruption (M2) rely exclusively
on self-reported survey measures rather than objective behavioral telemetry.
Consequently, we cannot verify exact AI utilization rates, such as API-call
volume or Copilot snippet-acceptance rates, limiting our claims to perceived
frequency and reported task orientation rather than technical dependency.
Consistent telemetry was not available across the proprietary IDE extensions,
hosted services, and student-owned development environments used in this
naturalistic bring-your-own-device setting. Multi-checkpoint self-reports were
therefore the only common measure of adoption that could be collected without
requiring access to vendor-controlled logs or intrusive monitoring of students'
devices.
Architectural Artifact Scope (M6a)
captures repository-visible T1 architectural artifacts, whereas Planning
Content Evidence (M6b) extracts structured evidence from commit subjects. The
Planning Content Evidence (M6b) remains exploratory and its categories were
subject to human review; neither measure establishes architectural quality or
planning intent. Clean-Path Rework (M8) uses file-path provenance to identify
changes to previously observed clean paths, not semantic defects or destructive
intent. Coordination-Language Evidence (M5) uses deterministic lexical markers
in transcript text; its counts are candidate
textual evidence, not direct diagnoses of coordination friction. These
distinctions constrain every metric's interpretation and prevent the metrics
from being treated as interchangeable measures of quality, effort, or
productivity.

The observational design and small sample constrain causal and external
validity. In particular, the late-stage concentration of Repository Activity
(M3) and Clean-Change Dynamics (M4) cannot be disentangled from ordinary
student procrastination (the ``student syndrome''). Without a non-AI control
group, we cannot determine whether these delays reflect GenAI-related
architectural debt or conventional late-stage integration pressure.
Artifact--Rework--Outcome Associations (M9) are exploratory, not causal or
confirmatory, and
leave-one-out ranges show sensitivity to individual team-semesters. Analytical
grains differ across student responses, team-semesters, rolling windows, and a
global transcript corpus, so denominators cannot be pooled implicitly. Rolling
windows overlap, so adjacent values are not independent observations. The
study covers 14 team-semesters in one educational setting; transcript evidence
is available only for 2025.2, which restricts cross-cohort comparability for
the lexical findings. Moreover, teams had complete autonomy over their process
models, including Scrum, Kanban, other Agile variants, and unguided approaches;
the study neither standardized these processes nor measured adherence to them.
We therefore cannot rule out that some integration friction arose from process
choice or poor methodological adherence rather than AI-tool use. Git cannot
observe planning or collaboration outside the repositories. Five Planning
Content Evidence (M6b) cases lacked T1
commit subjects and were
not measured; missing evidence was not converted into a low score. The
course required an integrated GenAI component but did not mandate
retrieval-augmented generation (RAG) or another common architecture; none of
the 14 projects implemented a full RAG pipeline. This removes mandatory RAG
complexity as a shared explanation for the observed integration friction.
Nevertheless, the teams' heterogeneous, self-selected AI integrations differed
in technical complexity, so architectural complexity remains a potential
project-level confound that this observational study cannot isolate. The
external SDD literature therefore motivates a testable intervention hypothesis,
but this study neither implements nor evaluates SDD. Its results support the
need to test more explicit specification and coordination practices under
Generative AI, not a claim that SDD has already improved outcomes.
